% Mục 4.4. So sánh đầu-cuối, bảng chính (khoảng 0,7 trang). Người viết: Thắng.
% Nguồn số: results/E-D1, results/E-D2/table.csv (code/experiments/e_d1.py, e_d2.py).
Khi chạy khối D trên khung và chữ chuẩn của tập kiểm tra (E-D1), F1 khoản thuốc đạt \num{1.000}.
Kết quả này chỉ cho thấy luật ghép nhất quán khi đầu vào đúng, chưa phải đánh giá độc lập, vì đáp án khoản thuốc cũng được dẫn xuất từ chính luật đó.
Kết quả chính của đồ án (E-D2) nằm ở Bảng~\ref{tab:ED2}.

\begin{table}[H]
\centering
\caption{Kết quả đầu-cuối trên 176 ảnh kiểm tra (E-D2). F1 khoản đòi hỏi tên đúng (CER $\le0{,}1$) và số lượng đúng; ngoặc vuông là khoảng tin cậy 95\% theo bootstrap.}\label{tab:ED2}
\footnotesize
\begin{tabularx}{\linewidth}{@{}Lcccccc@{}}
\toprule
\textbf{Pipeline} & \textbf{F1 khoản [CI 95\%]} & \textbf{Bỏ từ điển} & \textbf{F1 tên} & \textbf{SL đúng} & \textbf{Đơn đúng} & \textbf{s/ảnh} \\
\midrule
Chữ chuẩn (khung YOLO, chữ chuẩn)           & \num{1.0000} & \num{1.0000} & \num{1.0000} & \num{1.0000} & \num{1.0000} & --- \\
\midrule
(i) Cổ điển (B cổ điển, Tesseract, D)       & \num{0.9487} [\num{0.9316}; \num{0.9647}] & \num{0.9161} & \num{0.9580} & \num{0.9187} & \num{0.8011} & \textbf{\num{1.16}} \\
(ii) Lai (YOLO, Tesseract, D)              & \num{0.9865} [\num{0.9736}; \num{0.9977}] & \num{0.9549} & \num{0.9955} & \num{0.9865} & \num{0.9716} & \num{1.38} \\
(iv) YOLO, VietOCR công bố (C2), D          & \textbf{\num{0.9966}} [\num{0.9895}; \num{1.0000}] & \num{0.9514} & \textbf{\num{0.9966}} & \num{0.9955} & \textbf{\num{0.9943}} & \num{3.77} \\
(iii) YOLO, VietOCR tinh chỉnh (C3), D      & \textbf{\num{0.9966}} [\num{0.9930}; \num{1.0000}] & \textbf{\num{0.9966}} & \textbf{\num{0.9966}} & \textbf{\num{1.0000}} & \num{0.9830} & \num{3.37} \\
(v) VLM zero-shot (Qwen3.8-27B)              & \textbf{\num{0.9966}} [\num{0.9892}; \num{1.0000}] & \num{0.9875} & \textbf{\num{0.9966}} & \num{0.9932} & \textbf{\num{0.9943}} & --- \\
\bottomrule
\end{tabularx}
\par\smallskip{\footnotesize Nguồn: \path{result/e_d2_endtoend/}, \path{result/runtime/}. Cột ``Bỏ từ điển'' là F1 khoản khi tắt bước D4; cột ``s/ảnh'' là thời gian khối B+C (VLM gọi qua API từ xa, có 1/176 ảnh lỗi định dạng JSON).}
\end{table}

Pipeline chính (iii) đạt F1 khoản thuốc \num{0.9966} (CI 95\%: $[\num{0.9930}; \num{1.0000}]$), đúng 100\% số lượng (\num{1.0000}) và ngày kê đơn (\num{1.0000}), chỉ kém \num{0.0034} điểm so với mức trần dùng chữ chuẩn (\num{1.0000}).
Từ điển tên thuốc (D4) giúp tăng mạnh F1 khoản khi OCR còn lỗi ký tự: VietOCR công bố (iv) tăng \num{0.0452} điểm (từ \num{0.9514} lên \num{0.9966}), bản lai (ii) tăng \num{0.0316} điểm (\num{0.9549} lên \num{0.9865}), bản cổ điển (i) tăng \num{0.0326} điểm (\num{0.9161} lên \num{0.9487}) và VLM (v) tăng \num{0.0091} điểm (\num{0.9875} lên \num{0.9966}), trong khi VietOCR tinh chỉnh (iii) giữ nguyên \num{0.9966} cả khi bỏ từ điển nhờ CER tên thuốc chỉ còn \num{0.0003}.
So với VLM zero-shot Qwen3.8-27B (v), pipeline chính (iii) đạt cùng F1 khoản \num{0.9966} khi có từ điển nhưng cao hơn khi bỏ từ điển (\num{0.9966} so với \num{0.9875}), đạt độ chính xác tuyệt đối ở cả số lượng và ngày kê đơn (\num{1.0000} so với \num{0.9932} và \num{0.9943}), và không bị lỗi cú pháp đầu ra (VLM có 1/176 ảnh trả về JSON hỏng định dạng).
